\subsection{Examples of valuation rings}

\begin{enumerate}
\def\labelenumi{(\arabic{enumi})}
\item
  Every field is a valuation ring.
\item
  If \(\mathbf{V}'\) is a valuation ring of a field \(\mathbf{K}'\) and \(\mathbf{K}\) is a subfield of \(\mathbf{K}', \mathbf{V}' \cap \mathbf{K}\) is, by no. 2, Theorem 1 (c), a valuation ring of \(\mathbf{K}\) .
\item
  The following proposition provides numerous examples of valuation rings:
\end{enumerate}

PROPOSITION 2. Let A be a local ring whose maximal ideal is a principal ideal Ap. If \(\bigcap_{n=1}^{\infty} A p^n = (0)\) (for example if A is Noetherian, cf.~Chapter III, §3, no. 2,

Corollary to Proposition 5), the only ideals of \(\mathbf{A}\) are (0) and the \(\mathbf{A}p^n\) ; then either \(p\) is nilpotent or \(\mathbf{A}\) is a valuation ring.

Let A be filtered by the \(\mathbf{A}p^n\) and let \(v\) denote the corresponding order function (Chapter III, § 2, no. 2). As

\[
\bigcap_ {n = 1} ^ {\infty} A p ^ {n} = (0),
\]

the relation \(v(x) = +\infty\) implies \(x = 0\) . Let \(\mathfrak{a}\) be an ideal \(\neq (0)\) of \(\mathbf{A}\) and \(a\) an element of \(\mathfrak{a}\) at which \(v\) takes its least value; let us write \(v(a) = s\) ( \(s \neq +\infty\) ). Then \(\mathfrak{a} \subset \mathrm{Ap}^s\) . In particular, there exists \(u \in \mathbf{A}\) such that \(a = up^s\) ; as \(a \notin \mathrm{Ap}^{s+1}\) , \(u \notin \mathrm{Ap}\) ; hence \(u\) is invertible and \(p^s \in \mathbf{A}a \subset \mathfrak{a}\) . It follows that \(\mathfrak{a} = \mathrm{Ap}^s\) , whence our first assertion. It is also seen that every element \(a \neq 0\) of \(\mathbf{A}\) may be written in the form \(a = up^{v(a)}\) where \(u\) is invertible. If \(a' = u'p^{v(a')}\) ( \(u'\) invertible) is another non-zero element of \(\mathbf{A}\) , then \(aa' = uu'p^{v(a)+v(a')}\) ; hence, if \(p\) is not nilpotent, \(aa' \neq 0\) and \(\mathbf{A}\) is an integral domain. Then, as the set of ideals of \(\mathbf{A}\) is totally ordered by inclusion, we conclude that \(\mathbf{A}\) is a valuation ring (Theorem 1 (e)).

For example, if \(p\) is a prime number, the local ring \(\mathbf{Z}_{(p)}\) is a valuation ring. Let \(\mathbf{B} = \mathbf{K}[\mathbf{X}_1, \ldots, \mathbf{X}_n]\) be the polynomial ring in \(n\) indeterminates over a field \(\mathbf{K}\) ; the ideal \(\mathbf{BX}_1\) is prime, since \(\mathbf{B} / \mathbf{BX}_1\) is isomorphic to \(\mathbf{K}[\mathbf{X}_2, \ldots, \mathbf{X}_n]\) ; hence \(\mathbf{B}_{\mathbf{BX}_1}\) is a valuation ring; it is composed of rational functions \(\mathbf{PQ}^{-1}\) , where \(\mathbf{P}\) and \(\mathbf{Q}\) are polynomials and \(\mathbf{Q}(0, \mathbf{X}_2, \ldots, \mathbf{X}_n) \neq 0\) .

* More generally, we shall see that, if F is an extremal element of

\[
\mathrm{B} = \mathrm{K} [ \mathrm{X} _ {1}, \dots , \mathrm{X} _ {n} ],
\]

\(B_{BF}\) is a valuation ring (cf.~Chapter VII, § 3, no. 5). *

The ring of formal power series K{[}{[}T{]}{]} in one indeterminate over a field K is a Noetherian local integral domain whose maximal ideal is principal; it is therefore a valuation ring. On the other hand the ring K{[}{[}T₁, T₂{]}{]} of formal power series in two indeterminates, which is a Noetherian local integral domain is not a valuation ring, for neither of the elements T₁, T₂ is a multiple of the other.

PROPOSITION 3. Let A be a principal ideal domain and K its field of fractions. The valuation rings of K containing A and distinct from K are the rings of the form \(A_{Ap}\) , where p is an extremal element of A.

Clearly \(A_{Ap}\) (p extremal) is a valuation ring containing A and distinct from K (Proposition 2). Conversely, let V be a valuation ring distinct from K and containing A. As \(V \neq K\) , \(m(V)\) contains an element \(x \neq 0\) ; writing x = a/b where \(a \in A\) and \(b \in A\) , it is seen that \(A \cap m(V)\) contains the non-zero element a. As \(A \cap m(V)\) is prime, it is of the form Ap where p is extremal in A. Then \(A_{Ap} \subset V\) , \(pA_{Ap} \subset m(V)\) , so that V dominates \(A_{Ap}\) ; as \(A_{Ap}\) is a valuation ring of K (Proposition 2), \(V = A_{Ap}\) .

COROLLARY 1. Every valuation ring of the field \(\mathbf{Q}\) and distinct from \(\mathbf{Q}\) is of the form \(\mathbf{Z}_{(p)}\) where \(p\) is a prime number.

Every subring of Q contains Z.

COROLLARY 2. Let K be a field, K(X) the field of rational functions in one indeterminate over K and V a valuation ring of K(X) containing K and distinct from K(X). If X ∈ V, there exists an irreducible polynomial P ∈ K{[}X{]} such that \(V = (K{[}X{]})_{(P)}\); otherwise V is the local ring of K{[}X⁻¹{]} at the prime ideal X⁻¹K{[}X⁻¹{]} (in other words V is composed of the fractions A/B, where A ∈ K{[}X{]} and B ∈ K{[}X{]}, such that deg(A) ≤ deg(B)).

If \(\mathbf{X} \in \mathbf{V}\) , then \(\mathbf{K}[\mathbf{X}] \subset \mathbf{V}\) and the assertion made follows from Proposition 3. If \(\mathbf{X} \notin \mathbf{V}\) , then \(\mathbf{X}^{-1} \in \mathbf{V}\) , whence \(\mathbf{K}[\mathbf{X}^{-1}] \subset \mathbf{V}\) ; then \(\mathbf{V}\) is the local ring of \(\mathbf{K}[\mathbf{X}^{-1}]\) at a prime ideal \(p\) (Proposition 3) and this ideal contains \(\mathbf{X}^{-1}\) since \(\mathbf{X}^{-1}\) is not invertible in V; then \(p = X^{-1}K[X^{-1}]\) since the latter ideal is maximal. Finally let us consider a rational function \(A(X)/B(X)\) , where A and B are polynomials of respective degrees a and b; then \(A(X) = X^{a}A'(X^{-1})\) and \(B(X) = X^{b}B'(X^{-1})\) , where \(A'\) and \(B'\) are polynomials such that \(A'(0) \neq 0\) and \(B'(0) \neq 0\) ; hence, for \(A(X)/B(X)\) to belong to the local ring of \(K[X^{-1}]\) at \(X^{-1}K[X^{-1}]\) , it is necessary and sufficient that \(a \leqslant b\) .

